\section{Metodología CRISP-DM}

CRISP-DM organiza el trabajo en seis fases iterativas \citep{crispdm}. Un hallazgo de calidad,
modelamiento o evaluación puede obligar a revisar una fase anterior; el flujo no es estrictamente
lineal.

\begin{table}[H]
\centering
\caption{Estado de CRISP-DM en el proyecto.}
\label{tab:crispdm}
\begin{tabularx}{\textwidth}{p{3.2cm}X p{3cm}}
\toprule
Fase & Aplicación & Estado \\
\midrule
Business Understanding & Problema, interesados, objetivos, KPIs y restricciones & Completada
inicialmente \\
Data Understanding & Fuente, diccionario, calidad inicial y distribuciones & Completada inicialmente \\
Data Preparation & Índice, fila inválida, consolidación de IDs y géneros múltiples & Completada
inicialmente \\
Modeling & Líneas base y regresiones justificadas & Pendiente \\
Evaluation & MAE, RMSE, $R^2$, utilidad, errores por grupos y robustez temporal & Pendiente \\
Deployment & Repositorio reproducible, informe y presentación & En construcción \\
\bottomrule
\end{tabularx}
\end{table}

La separación entre notebooks refleja responsabilidades: \code{01_data_understanding.ipynb}
documenta negocio y diagnóstico; \code{02_data_preparation.ipynb} transforma la fuente;
\code{03_exploratory_analysis.ipynb} estudia canciones únicas. La lógica común se implementa en
\code{src/spotify_popularity/} y se verifica en \code{tests/}.
